% =====================================================================
\chapter{La descomposición en valores singulares}
\label{tema:4}
% =====================================================================

El teorema espectral solo se aplica a matrices cuadradas simétricas. La descomposición en valores singulares (SVD) extiende la misma idea a cualquier matriz rectangular: toda transformación lineal es un giro, seguido de un estiramiento a lo largo de ejes perpendiculares, seguido de otro giro. Este tema sigue a Strang~\cite[§7.1]{strang-ila}, la clase~29 de 18.06~\cite{mit1806}, las clases~6 y~8 de 18.065~\cite{mit18065} y el §4.5 de~\cite{mml}.

\section{El teorema de la SVD}

\begin{teorema}[Descomposición en valores singulares]\label{teo:svd}
Sea $A\in\R^{m\times n}$ de rango $r$. Existen matrices ortogonales $U\in\R^{m\times m}$ y $V\in\R^{n\times n}$ y números $\sigma_1\ge\sigma_2\ge\dots\ge\sigma_r>0$ tales que
\[
A=U\Sigma V\T=\sum_{i=1}^{r}\sigma_i\,u_iv_i\T,
\]
donde $\Sigma\in\R^{m\times n}$ tiene $\Sigma_{ii}=\sigma_i$ para $i\le r$ y ceros en el resto, y $u_i$, $v_i$ son las columnas de $U$ y $V$. Los $\sigma_i$ son los \emph{valores singulares} de $A$, y los $u_i$ y $v_i$, sus \emph{vectores singulares} izquierdos y derechos. Además, $Av_i=\sigma_iu_i$ y $A\T u_i=\sigma_iv_i$ para $i\le r$.
\end{teorema}

\addtocounter{figure}{1}%

\begin{observacion}[SVD reducida]
Como las columnas $u_i$ y $v_i$ con $i>r$ se multiplican por ceros, también se cumple $A=U_r\Sigma_rV_r\T$, con $U_r\in\R^{m\times r}$, $\Sigma_r=\diag(\sigma_1,\dots,\sigma_r)$ y $V_r\in\R^{n\times r}$. Es la versión que se almacena en la práctica (\texttt{np.linalg.svd} con \texttt{full\_matrices=False}).
\end{observacion}

\begin{corolario}\label{cor:svd-props}
Sea $A=U\Sigma V\T$ la SVD de $A\in\R^{m\times n}$.
\begin{enumerate}
  \item $\rango(A)$ es el número de valores singulares no nulos.
  \item $\sigma_i(A)=\sqrt{\lambda_i(A\T A)}$, y $AA\T$ tiene los mismos autovalores no nulos que $A\T A$, con autovectores $u_i$.
  \item $\{u_1,\dots,u_r\}$ es una base ortonormal de $\Col(A)$ y $\{v_{r+1},\dots,v_n\}$, una base ortonormal de $\Nuc(A)$.
  \item Si $A$ es simétrica con autovalores $\lambda_i$, sus valores singulares son los $\abs{\lambda_i}$. Si además es semidefinida positiva, la SVD coincide con la descomposición espectral.
\end{enumerate}
\end{corolario}

\section{Cálculo a mano}

El procedimiento de la demostración sirve también para calcular la SVD de matrices pequeñas: (1) se forma $A\T A$ o $AA\T$, la que sea de menor orden; (2) se calculan sus autovalores $\sigma_i^{2}$ y autovectores ortonormales; (3) se obtienen los vectores singulares del otro lado con $u_i=Av_i/\sigma_i$ o $v_i=A\T u_i/\sigma_i$.

\addtocounter{definicion}{1}%

\section{Normas matriciales}

\begin{definicion}[Normas de Frobenius y espectral]
Para $A\in\R^{m\times n}$ se definen la \emph{norma de Frobenius} y la \emph{norma espectral}
\[
\norm{A}_F=\Bigl(\sum_{i=1}^{m}\sum_{j=1}^{n}a_{ij}^{2}\Bigr)^{1/2},\qquad
\norm{A}_2=\max_{x\ne 0}\frac{\norm{Ax}}{\norm{x}}.
\]
\end{definicion}

\begin{proposicion}\label{prop:normas-svd}
Para toda $A\in\R^{m\times n}$ con valores singulares $\sigma_1\ge\dots\ge\sigma_r>0$,
\[
\norm{A}_F^{2}=\traza(A\T A)=\sum_{i=1}^{r}\sigma_i^{2},\qquad \norm{A}_2=\sigma_1.
\]
En particular, $\norm{A}_2\le\norm{A}_F\le\sqrt{r}\,\norm{A}_2$.
\end{proposicion}

Las dos igualdades se demuestran en los \cref{eje:4-frobenius,eje:4-espectral}; la segunda es el teorema~4.24 de~\cite{mml}. Ambas normas son invariantes por transformaciones ortogonales: $\norm{Q_1AQ_2}=\norm{A}$ para $Q_1$, $Q_2$ ortogonales, porque los valores singulares no cambian. Las desigualdades finales se siguen de $\sigma_1^{2}\le\sum_i\sigma_i^{2}\le r\sigma_1^{2}$.

\addtocounter{definicion}{1}%

\addtocounter{definicion}{1}%

